\documentclass[12pt, a4paper]{article}

% Поддержка русского языка и кодировки
\usepackage[utf8]{inputenc}
\usepackage[T2A]{fontenc}
\usepackage[russian]{babel}

% Настройка полей
\usepackage[left=2.5cm, right=1.5cm, top=2cm, bottom=2cm]{geometry}

% Математические пакеты
\usepackage{amsmath, amsfonts, amssymb}

% Пакет для работы с графикой
\usepackage{graphicx}

% Настройка абзацного отступа и интерлиньяжа
\usepackage{indentfirst}
\linespread{1.25}

% Пакет для кликабельного оглавления и ссылок
\usepackage{hyperref}
\hypersetup{
    colorlinks=true,
    linkcolor=black,
    filecolor=magenta,      
    urlcolor=blue,
}

\usepackage{listings}
\usepackage{color}
\usepackage{float}

\definecolor{codegreen}{rgb}{0,0.6,0}
\definecolor{codeorange}{rgb}{0.8, 0.5, 0}

\lstdefinestyle{customc}{
  inputencoding=utf8x,
  extendedchars=\true,
  belowcaptionskip=1\baselineskip,
  breaklines=true,
  frame=L,
  xleftmargin=\parindent,
  language=C++,
  showstringspaces=false,
  basicstyle=\footnotesize\ttfamily,
  keywordstyle=\bfseries\color{blue},
  commentstyle=\color{codegreen},
  identifierstyle=\color{black},
  stringstyle=\color{codeorange},
  keepspaces=true,
}
\lstset{style=customc}

\begin{document}

% --- ТИТУЛЬНЫЙ ЛИСТ ---
\begin{titlepage}

\begin{center}
Министерство науки и высшего образования Российской Федерации \\
Федеральное государственное автономное образовательное учреждение \\
высшего образования\\ 
\textbf{«Национальный исследовательский \\ Нижегородский государственный университет \\
им. Н.И. Лобачевского (ННГУ)»}\\[1.5cm]
\textbf{Институт информационных технологий, математики и механики }\\[5.5cm]

\textbf{Отчет по лабораторной работе} \\
\textbf{по курсу}  \\
``Параллельные численные методы'' \\
``Метод сопряженных градиентов'' \\ [5cm]


\begin{flushright}
\begin{minipage}{0.52\textwidth}
\begin{flushleft}
\textbf{Выполнил:} \\
студент группы 3825М1ПМкн1 \\
Сучков В.Н.\\[0.3cm]
\textbf{Проверил:}\\
д.т.н., проф. кафедры МОСТ\\  Баркалов К.А.\\[0.3cm]
\end{flushleft}
\end{minipage}
\end{flushright}

\vfill

Нижний Новгород \\
2025

\thispagestyle{empty}
\end{center}

\newpage
\end{titlepage}

% --- ОГЛАВЛЕНИЕ ---
\newpage
\tableofcontents
\newpage    

% --- ВВЕДЕНИЕ ---
\section{Введение}
Целью данной работы является реализация, тестирование и исследование производительности 
параллельного алгоритма метода сопряженных градиентов для решения 
СЛАУ вида $Ax = b$, где $A = A^T > 0$.

В современных вычислительных задачах (например, при решении дифференциальных уравнений в частных производных сеточными методами) 
матрицы систем достигают огромных размеров, однако подавляющее большинство их элементов равны нулю. 
Хранение таких матриц в плотном формате нецелесообразно из-за колоссальных затрат памяти и 
выполнения множества избыточных операций умножения на ноль. В связи с этим в работе применяется формат сжатого хранения CRS. 
Для дополнительной оптимизации, учитывая симметрию матрицы, 
в памяти хранятся только элементы, расположенные не ниже главной диагонали.

Метод сопряженных градиентов относится к классу итерационных методов и является одним из наиболее эффективных способов решения СЛАУ с симметричными положительно определенными матрицами. 
Основной вычислительной нагрузкой алгоритма является операция умножения разреженной матрицы на плотный вектор. 
Для ускорения расчетов и распараллеливания вычислений в системах с общей памятью применяется фреймворк OpenMP. 

Особую сложность при распараллеливании представляет структура симметричного формата CRS: при умножении матрицы на вектор возникает необходимость одновременной записи результатов вычислений симметричных 
внедиагональных элементов разными потоками в одни и те же ячейки результирующего вектора. Устранение возникающего состояния гонки данных и эффективное распределение нагрузки между вычислительными 
ядрами является ключевой задачей данной реализации.

\section{Постановка задачи}
Требуется реализовать на языке C++ с использованием технологии OpenMP метод сопряженных градиентов для решения СЛАУ вида:
\begin{equation}
    Ax = b
\end{equation}
где $A = A^T > 0$ — разреженная матрица размерности $n \times n$, а $x$ и $b$ — плотные векторы размера $n \times 1$.

Программа должна реализовывать функцию со следующим заголовком:
\begin{lstlisting}[language=C++, frame=single]
void SLE_Solver_CRS(CRSMatrix & A, double * b, double eps, 
                    int max_iter, double * x, int & count);
\end{lstlisting}

\subsection{Формат входных и выходных данных}
\begin{itemize}
    \item \textbf{Входные параметры:}
    \begin{itemize}
        \item \texttt{A} — ссылка на структуру \texttt{CRSMatrix}, содержащую матрицу в симметричном формате CRS (хранятся только элементы не ниже главной диагонали, т.е. $j \ge i$).
        \item \texttt{b} — указатель на массив правой части СЛАУ размера $n$.
        \item \texttt{eps} — критерий остановки итерационного процесса по относительной норме изменения решения.
        \item \texttt{max\_iter} — максимальное число итераций.
    \end{itemize}
    \item \textbf{Выходные параметры:}
    \begin{itemize}
        \item \texttt{x} — указатель на массив для записи результирующего вектора решения.
        \item \texttt{count} — количество фактически выполненных алгоритмом итераций.
    \end{itemize}
\end{itemize}

Матрица $A$ представляется специализированной структурой:
\begin{lstlisting}[language=C++, frame=single]
struct CRSMatrix {
    int n;           // Число строк в матрице
    int m;           // Число столбцов в матрице
    int nz;          // Число ненулевых элементов в верхнем треугольнике
    vector<double> val;       // Массив значений элементов (по строкам)
    vector<int> colIndex;     // Массив номеров столбцов
    vector<int> rowPtr;       // Массив индексов начала строк
};
\end{lstlisting}

\subsection{Критерии успешности и ограничения}
\begin{itemize}
    \item \textbf{Ограничения на размер задачи:} размерность матрицы $n \le 100000$, общее число ненулевых элементов $nz \le 10^7$. Лимит по времени: 100 с. Лимит по памяти: 256 МБ.
    \item \textbf{Критерий остановки внутри цикла:} 
    \begin{equation}
        \frac{\|x_k - x_{k+1}\|_2}{\|b\|_2} < \text{eps}
    \end{equation}
    или превышение лимита итераций \texttt{max\_iter}.
    \item \textbf{Критерий корректности ответа:} Относительная норма невязки на выходе должна удовлетворять условию:
    \begin{equation}
        \frac{\|Ax - b\|_2}{\|A\|_2} < 10^{-8}
    \end{equation}
    \item \textbf{Требование к масштабируемости:} Эффективность параллельной версии должна составлять не менее 50\% на больших тестах.
\end{itemize}

\pagebreak
\section{Метод решения задачи}
Для решения СЛАУ применяется итерационный метод сопряженных градиентов. Метод строит последовательность приближений $x_k$ и сходится за конечное число шагов.

\subsection{Математический алгоритм}
Пусть задано начальное приближение $x_0 = 0$. Тогда инициализация алгоритма имеет вид:
\begin{align*}
    r_0 &= b - Ax_0 = b \\
    p_0 &= r_0
\end{align*}
На каждой итерации $k = 0, 1, \dots, \text{max\_iter}-1$ выполняются следующие вычисления:
\begin{enumerate}
    \item Вычисление матрично-векторного произведения: $A_p = A \cdot p_k$.
    \item Вычисление шага вдоль направления:
    \begin{equation}
        \alpha_k = \frac{(r_k, r_k)}{(A_p, p_k)}
    \end{equation}
    \item Обновление вектора решения: $x_{k+1} = x_k + \alpha_k p_k$.
    \item Проверка критерия остановки по норме приращения решения.
    \item Обновление невязки: $r_{k+1} = r_k - \alpha_k A_p$.
    \item Вычисление коэффициента сопряженности:
    \begin{equation}
        \beta_k = \frac{(r_{k+1}, r_{k+1})}{(r_k, r_k)}
    \end{equation}
    \item


Построение нового направления: $p_{k+1} = r_{k+1} + \beta_k p_k$.
\end{enumerate}

\subsection{Особенности умножения в симметричном CRS}
Наиболее трудоемкой операцией метода является произведение разреженной матрицы на веткор 
$A \cdot p_k$. Так как в структуре \texttt{CRSMatrix} хранятся только элементы верхнего треугольника ($j \ge i$), 
стандартный обход строк не позволит вычислить полное произведение. 

Для корректного умножения каждый элемент $v = A_{ij}$ ($i \neq j$) обрабатывается дважды:
\begin{itemize}
    \item Прямой вклад в строку $i$: $\text{sum} \mathrel{+}= v \cdot p_j$
    \item Симметричный вклад в строку $j$: $A_p[j] \mathrel{+}= v \cdot p_i$
\end{itemize}
Такой подход минимизирует расходы памяти, но порождает серьезную проблему для параллельной реализации, 
так как разные потоки могут одновременно пытаться обновить одну и ту же ячейку массива $A_p[j]$. Описание борьбы с этой проблемой приведено в следующих разделах.

% --- ОПИСАНИЕ РЕАЛИЗАЦИИ АЛГОРИТМА ---
\section{Описание реализации алгоритма}

\subsection{Используемые структуры данных}
Для эффективной работы с разреженными матрицами больших размерностей в программе используется симметричный формат Compressed Row Storage. В отличие от классического плотного представления, данный формат позволяет хранить исключительно ненулевые элементы структуры, расположенные на главной диагонали и выше нее (в силу симметрии матрицы $A$). 

Память под матрицу выделяется динамически в рамках структуры \texttt{CRSMatrix}:
\begin{itemize}
    \item \texttt{val} — одномерный массив вещественных чисел типа \texttt{double}, содержащий значения ненулевых элементов верхнего треугольника матрицы, записанные последовательно по строкам.
    \item \texttt{colIndex} — одномерный массив целых чисел типа \texttt{int}, хранящий номера столбцов для каждого соответствующего элемента из массива \texttt{val}.
    \item \texttt{rowPtr} — одномерный массив целых чисел типа \texttt{int} размера $n+1$, где элемент с индексом $i$ указывает на позицию начала $i$-й строки в массивах \texttt{val} и \texttt{colIndex}. Последний элемент \texttt{rowPtr[n]} равен общему числу ненулевых элементов верхнего треугольника (\texttt{nz}).
\end{itemize}

Вспомогательные векторы $r$ (текущяя невязка), $p$ (направление), $A_p$ (результат умножения матрицы на вектор направления), 
а также вектор искомого решения $x$ реализуются стандартными динамическими массивами \texttt{std::vector<double>} и располагаются в памяти непрерывно.

\subsection{Описание базовой логики программы}
Реализация алгоритма сосредоточена в функции \texttt{SLE\_Solver\_CRS} и разбита на три ключевых этапа:
\begin{enumerate}
    \item \textbf{Инициализация и выделение ресурсов:} 
    Начальное приближение $x_0$ заполняется нулями. Вектор остатка $r$ и начальное направление спуска $p$ инициализируются значениями вектора правой части $b$. Вычисляется квадрат нормы вектора правой части $\text{norm\_b\_sq} = (b, b)$ и его скалярный корень для последующих проверок сходимости. Если норма равна нулю, программа мгновенно завершает работу, возвращая нулевой вектор.
    
    \item \textbf{Выделение локальных буферов потоков:}
    Для предотвращения конфликтов параллельного доступа создается двумерный массив \texttt{local\_Ap} размера $\text{max\_threads} \times n$. Каждому выполняющемуся потоку сопоставляется собственный изолированный вектор для накопления частичных результатов матрично-векторного умножения.
    
    \item \textbf{Основной итерационный цикл:}
    Цикл выполняется до достижения \texttt{max\_iter} или досрочного выполнения критерия остановки. Каждая итерация включает в себя:
    \begin{itemize}
        \item Обнуление локальных буферов \texttt{local\_Ap} для текущего шага.
        \item Вычисление произведения $A \cdot p$ с учетом симметрии матрицы.
        \item Редукцию (суммирование) локальных буферов всех потоков в глобальный вектор $A_p$.
        \item Вычисление скалярного произведения $(A_p, p_k)$ и корректировку шага $\alpha_k$.
        \item Модификацию решения $x$ и вычисление нормы его приращения $\Delta x$.
        \item Обновление вектора невязки $r$, расчет коэффициента $\beta_k$ и пересчет направления $p$.
    \end{itemize}
\end{enumerate}

% --- РАЗЛИЧНЫЕ СПОСОБЫ РАСПАРАЛЛЕЛИВАНИЯ И ИХ СРАВНЕНИЕ ---
\section{Различные способы распараллеливания и их сравнение}

\subsection{Варианты декомпозиции данных}
При распараллеливании операции умножения разреженной симметричной матрицы на плотный вектор ($A \cdot p$) 
по строкам возникает проблема гонки данных. 

Рассмотрим два альтернативных подхода к организации вычислений:
\begin{enumerate}
    \item \textbf{Прямое распараллеливание с атомарными операциями:}
    Потоки распределяют итерации внешнего цикла по строкам $i$. Внутри цикла элемент $A_{ij}$ умножается на $p_j$ и добавляется в локальную сумму строки $i$. Симметричный вклад

$A_{ij} \cdot p_i$ должен быть добавлен в ячейку $A_p[j]$. Поскольку разные строки (и, соответственно, разные потоки) могут иметь элементы, относящиеся к одному и тому же столбцу $j$, запись в $A_p[j]$ требует синхронизации (например, через директиву \texttt{\#pragma omp atomic}). Эксперименты показывают, что высокая плотность атомарных операций катастрофически снижает производительность из-за блокировок шины данных.
    
    \item \textbf{Декомпозиция с использованием приватных буферов (выбранный подход):}
    Каждому потоку выделяется индивидуальный вектор \texttt{local\_Ap[tid]} размера $n$. Потоки полностью изолированы друг от друга в процессе обхода матрицы. Вклад от внедиагонального элемента записывается потоком в свой собственный буфер: \texttt{local\_Ap[tid][j] += v * p[i]}. Гонка данных полностью исключается на этапе расчета, а накладные расходы сводятся к финальной параллельной сборке (редукции) буферов, выполняемой за $O(\text{n\_threads} \cdot n)$ операций.
\end{enumerate}

% --- РЕАЛИЗАЦИЯ ПАРАЛЛЕЛЬНОГО АЛГОРИТМА С ИСПОЛЬЗОВАНИЕМ ОБЩЕЙ ПАМЯТИ ---
\section{Реализация параллельного алгоритма с использованием общей памяти}

\subsection{Организация параллельных секций}
Для минимизации накладных расходов на повторное создание и уничтожение потоков, основной 
вычислительный цикл по итерациям метода сопряженных градиентов помещен внутрь единой параллельной области 
\texttt{\#pragma omp parallel}.

\begin{lstlisting}[language=C++, frame=single]
#pragma omp parallel
{
    int tid = omp_get_thread_num();
    int n_threads = omp_get_num_threads();
    // Дальнейший итерационный цикл по шагам метода
}
\end{lstlisting}

\subsection{Синхронизация и последовательные участки}
Внутри основной параллельной области явные барьеры синхронизации отсутствуют, поскольку логика алгоритма опирается на неявные барьеры, встроенные в директивы распределения циклов \texttt{\#pragma omp for}. 

Потоки синхронизируются в следующих критических точках:
\begin{itemize}
    \item \textbf{После обнуления локальных буферов:} Перед началом умножения матрицы на вектор гарантируется, что все потоки завершили очистку своих строк в \texttt{local\_Ap}.
    \item \textbf{После этапа умножения разреженной матрцы на вектор:} Перед началом финального суммирования все частичные вклады потоков должны быть полностью вычислены.
    \item \textbf{После сборки глобального вектора $A_p$:} Неявный барьер в конце цикла сборки гарантирует, что к моменту перехода к вычислению скалярного произведения $(A_p, p)$ вектор $A_p$ полностью сформирован.
\end{itemize}

\subsection{Параллельная обработка разреженной структуры}
Основная вычислительная фаза умножения симметричной разреженной матрицы на вектор реализована следующим образом:

\begin{lstlisting}[language=C++, frame=single]
#pragma omp for
for (int i = 0; i < n; i++) {
    double sum = 0.0;
    for (int k = A.rowPtr[i]; k < A.rowPtr[i+1]; k++) {
        int j = A.colIndex[k];
        double v = A.val[k];
        
        sum += v * p[j]; // Вклад в текущую строку i
        
        if (i != j) {
            local_Ap[tid][j] += v * p[i]; % Симметричный вклад в строку j
        }
    }
    local_Ap[tid][i] += sum;
}
\end{lstlisting}

Распределение итераций внешнего цикла по строкам $i$ выполняется автоматически между доступными потоками. Каждый поток накапливает результаты в свою изолированную строку \texttt{local\_Ap[tid]}, что исключает состояние гонки при обработке внедиагональных элементов.


% --- ОПИСАНИЕ ТЕСТОВЫХ ЗАДАЧ ---
\section{Описание тестовых задач}
Для всестороннего анализа корректности и эффективности разработанного параллельного алгоритма было проведено комплексное тестирование, разделенное на два этапа.

\subsection{Тесты на проверку корректности}
Целью данного этапа является верификация математической точности реализованного алгоритма и проверка соблюдения критерия сходимости. Для проведения экспериментов генерировались случайные разреженные квадратные матрицы со строгим диагональным преобладанием, что гарантирует их симметричность и положительную определенность.

Тестирование проводилось на наборах матриц различных размерностей:
\begin{itemize}
    \item \textbf{Малые разреженные матрицы ($n \le 1000$):} Использовались для базовой проверки логики итерационного процесса, корректности индексации массивов CRS (\texttt{rowPtr}, \texttt{colIndex}) и обработки граничных случаев (например, строк, содержащих только диагональный элемент).
    \item \textbf{Средние разреженные матрицы ($n = 10000$):} Применялись для подтверждения устойчивости алгоритма при большом числе итераций и проверки корректности работы критерия остановки \texttt{eps}.
\end{itemize}
Ответ считался корректным, если относительная норма остатка удовлетворяла условию:
\begin{equation}
    \frac{\|Ax - b\|_2}{\|A\|_2} < 10^{-8}
\end{equation}

\subsection{Тесты производительности и масштабируемости}
Для исследования скоростных характеристик


параллельной версии требовалось создать существенную вычислительную нагрузку, минимизирующую долю системных накладных расходов OpenMP на создание и синхронизацию потоков. 

В соответствии с техническими требованиями, ограничения на размер задачи составляют $n \le 100000$ при количестве ненулевых элементов $nz \le 10^7$. Для профилирования масштабируемости была выбрана тестовая конфигурация со следующими параметрами:
\begin{itemize}
    \item Размерность матрицы: $n = 100000$.
    \item Число ненулевых элементов: $nz = 10^7$.
    \item Диапазон изменения значений векторов: $|x_i| \le 1, i = \overline{1, n}$.
\end{itemize}

Замеры общего времени выполнения алгоритма производились для последовательного режима (1 поток), 
а также для параллельного выполнения на конфигурациях с числом потоков $p \in \{2, 3, 4, 5, 6\}$. 
На основе полученных временных данных рассчитывались коэффициенты ускорения $S_p$ и эффективности $E_p$,
 что позволяет оценить степень масштабируемости алгоритма и выявить точку насыщения пропускной способности 
 памяти вычислительной системы.

% --- РЕЗУЛЬТАТЫ ЭКСПЕРИМЕНТОВ ---
\section{Результаты экспериментов}

\subsection{Проверка корректности}
Для верификации работы алгоритма проводилось вычисление относительной невязки полученного решения. В таблице 1 представлены результаты запуска на сгенерированных случайных матрицах с диагональным преобладанием.

\begin{table}[h]
\centering
\begin{tabular}{|c|c|c|}
\hline
\textbf{Размер матрицы ($n$)} & \textbf{Итераций} & \textbf{Относительная невязка}  \\ 
\hline
10 & 11 & $8.524502 \times 10^{-17}$  \\
50 & 22 & $1.342407 \times 10^{-9}$ \\
100 & 24 & $1.235791 \cdot 10^{-9}$  \\
500 & 27 & $5.558743 \cdot 10^{-9}$ \\
1000 & 29 & $9.014339 \cdot 10^{-9}$  \\
\hline
\end{tabular}
\caption{Результаты проверки корректности алгоритма}
\end{table}

Как видно из таблицы, на всех тестовых размерностях алгоритм сходится, а относительная невязка строго удовлетворяет требуемому условию $\frac{\|Ax - b\|_2}{\|A\|_2} < 10^{-8}$. Небольшой рост невязки с увеличением размерности является математически ожидаемым и объясняется стандартным накоплением вычислительных погрешностей округления.

\subsection{Тесты производительности и масштабируемости}
Для исследования масштабируемости параллельной версии использовалась матрица размерности $n = [N]$, что позволило создать достаточную вычислительную нагрузку и нивелировать влияние накладных расходов OpenMP.

В таблице 2 представлены результаты замеров времени выполнения, а также рассчитанные значения ускорения ($S = T_1 / T_p$) и эффективности ($E = S / p$).

\begin{table}[h]
\centering
\begin{tabular}{|c|c|c|c|}
\hline
\textbf{Число потоков ($p$)} & \textbf{Время, с} & \textbf{Ускорение ($S$)} & \textbf{Эффективность ($E$), \%} \\ \hline
1 & 1.906 & 1.00 & 100.0 \\
2 & 1.277 & 1.49 & 74.63 \\
3 &  0.958 &  1.99 & 66.32 \\
4 &  0.787  & 2.42 & 60.55 \\
5 &  0.6810   &  2.8  & 55.98 \\
6&  0.621  &  3.07   & 51.15 \\
\hline
\end{tabular}
\caption{Производительность при $n=100000$}
\end{table}

Анализ полученных данных показывает уверенную масштабируемость алгоритма на выбранной архитектуре. Эффективность распараллеливания остается на высоком уровне и на всех тестовых конфигурациях строго превышает заданный техническим заданием порог в 50\%. Использование локальных буферов позволило избежать блокировок и обеспечить равномерную загрузку вычислительных ядер.

% --- ЗАКЛЮЧЕНИЕ ---
\section{Заключение}
В рамках лабораторной работы была успешно реализована, протестирована и проанализирована параллельная 
версия метода сопряженных градиентов для систем линейных алгебраических уравнений 
с разреженными матрицами с использованием фреймворка OpenMP.

Применение симметричного формата CRS позволило существенно оптимизировать потребление 
оперативной памяти. Ключевая проблема распараллеливания — состояние гонки данных
при записи симметричных вкладов — была эффективно решена за счет 
выделения приватных потоковых массивов памяти. Это избавило код от 
необходимости использования ресурсоемких атомарных операций в основном вычислительном цикле.

По итогам тестирования подтверждено:
\begin{itemize}
    \item \textbf{Корректность:} алгоритм стабильно сходится к правильному решению, выполняя жесткие 
    критерии по относительной норме невязки ($< 10^{-8}$) на всех тестовых размерностях.
    \item \textbf{Производительность:} параллельная версия алгоритма демонстрирует высокую масштабируемость. 
    Требование к сохранению эффективности распараллеливания на уровне не менее 50\% выполнено в 
    полном объеме для всех протестированных конфигураций потоков.
\end{itemize}

Поставленная задача решена полностью, заявленные ограничения по памяти и времени соблюдены.

% --- СПИСОК ЛИТЕРАТУРЫ ---
\section{Список литературы}
\begin{enumerate}
    \item Баркалов К.А. Параллельные численные методы: презентации лекций. ННГУ им. Н.И. Лобачевского.
    \item Самарский А.А., Гулин А.В. Численные методы. М.: Наука, 1989.
    \item Антонов А.С. Технологии параллельного программирования MPI и OpenMP: Учебное пособие. — М.: Издательство
\end{enumerate}

\pagebreak

\appendix
\section{Приложение. Код}

\begin{lstlisting}[language=C++, frame=single]
#include <iostream>
#include <vector>
#include <cmath>
#include <omp.h>
#include <random>
#include <iomanip>
#include <chrono>

struct CRSMatrix {
    int n; 
    int m; 
    int nz; 
    std::vector<double> val;
    std::vector<int> colIndex;
    std::vector<int> rowPtr;
};


void SLE_Solver_CRS(CRSMatrix & A, double * b, double eps, int max_iter, double * x, int & count)
{
    int n = A.n;
    count = 0;

    std::vector<double> r(n, 0.0);
    std::vector<double> p(n, 0.0);
    std::vector<double> Ap(n, 0.0);

    int max_threads = omp_get_max_threads();
    std::vector<std::vector<double>> local_Ap(max_threads, std::vector<double>(n, 0.0));

    double norm_b_sq = 0.0;

    #pragma omp parallel for reduction(+:norm_b_sq)
    for (int i = 0; i < n; i++) {
        x[i] = 0.0;
        r[i] = b[i];
        p[i] = b[i];
        norm_b_sq += b[i] * b[i];
    }

    double norm_b = sqrt(norm_b_sq);
    
    if (norm_b == 0.0) {
        return;
    }

    double rr = norm_b_sq;

    for (int iter = 0; iter < max_iter; iter++) {
        
        #pragma omp parallel
        {
            int tid = omp_get_thread_num();
            int n_threads = omp_get_num_threads();

            for (int i = 0; i < n; i++) {
                local_Ap[tid][i] = 0.0;
            }

            #pragma omp for
            for (int i = 0; i < n; i++) {
                double sum = 0.0;
                for (int k = A.rowPtr[i]; k < A.rowPtr[i+1]; k++) {
                    int j = A.colIndex[k];
                    double v = A.val[k];
                    
                    sum += v * p[j];
                    
                    if (i != j) {
                        local_Ap[tid][j] += v * p[i];
                    }
                }
                local_Ap[tid][i] += sum;
            } 

            #pragma omp for
            for (int i = 0; i < n; i++) {
                double total = 0.0;
                for (int t = 0; t < n_threads; t++) {
                    total += local_Ap[t][i];
                }
                Ap[i] = total;
            }
        } 

        double pAp = 0.0;
        #pragma omp parallel for reduction(+:pAp)
        for (int i = 0; i < n; i++) {
            pAp += p[i] * Ap[i];
        }

        double alpha = rr / pAp;

        double dx_norm_sq = 0.0;
        #pragma omp parallel for reduction(+:dx_norm_sq)
        for (int i = 0; i < n; i++) {
            double dx = alpha * p[i];
            x[i] += dx;
            dx_norm_sq += dx * dx;
        }

        count++;

        double dx_norm = sqrt(dx_norm_sq);
        if (dx_norm / norm_b < eps) {
            break;
        }

        double rr_new = 0.0;
        #pragma omp parallel for reduction(+:rr_new)
        for (int i = 0; i < n; i++) {
            r[i] -= alpha * Ap[i];
            rr_new += r[i] * r[i];
        }

        double beta = rr_new / rr;

        #pragma omp parallel for
        for (int i = 0; i < n; i++) {
            p[i] = r[i] + beta * p[i];
        }

        rr = rr_new;
    }
}

void multiply_CRS_vector(const CRSMatrix& A, const std::vector<double>& x, std::vector<double>& res) {
    int n = A.n;
    fill(res.begin(), res.end(), 0.0);
    for (int i = 0; i < n; i++) {
        double sum = 0.0;
        for (int k = A.rowPtr[i]; k < A.rowPtr[i+1]; k++) {
            int j = A.colIndex[k];
            double v = A.val[k];
            sum += v * x[j];
            if (i != j) {
                res[j] += v * x[i];
            }
        }
        res[i] += sum;
    }
}


CRSMatrix generate_spd_crs(int n, int max_nz_per_row) {
    CRSMatrix A;
    A.n = n;
    A.m = n;
    A.rowPtr.push_back(0);
    A.nz = 0;

    std::mt19937 gen(144); 
    std::uniform_real_distribution<double> dist_val(-10.0, 10.0);
    std::uniform_int_distribution<int> dist_col(0, n - 1);

    std::vector<double> diag_vals(n, 0.0);
    std::vector<std::vector<std::pair<int, double>>> rows(n);

    for (int i = 0; i < n; i++) {
        rows[i].push_back({i, 0.0}); 
        
        int nz_count = dist_col(gen) % max_nz_per_row;
        for (int k = 0; k < nz_count; k++) {
            int j = dist_col(gen);
            if (j > i) {
                double val = dist_val(gen);
                rows[i].push_back({j, val});
                diag_vals[i] += abs(val);
                diag_vals[j] += abs(val);
            }
        }
    }

    for (int i = 0; i < n; i++) {
        rows[i][0].second = diag_vals[i] + 10.0 + abs(dist_val(gen));
        
        for (auto& elem : rows[i]) {
            A.colIndex.push_back(elem.first);
            A.val.push_back(elem.second);
            A.nz++;
        }
        A.rowPtr.push_back(A.nz);
    }
    return A;
}

double calc_matrix_2norm(const CRSMatrix& A) {
    int n = A.n;
    std::vector<double> q(n, 1.0 / sqrt(n));
    std::vector<double> z(n, 0.0);
    double lambda = 0.0;

    for (int iter = 0; iter < 100; iter++) {
        multiply_CRS_vector(A, q, z);
        double z_norm = 0.0;
        for (int i = 0; i < n; i++) z_norm += z[i] * z[i];
        z_norm = sqrt(z_norm);
        
        double new_lambda = 0.0;
        for (int i = 0; i < n; i++) {
            q[i] = z[i] / z_norm;
            new_lambda += q[i] * z[i]; // (Aq, q)
        }
        if (abs(new_lambda - lambda) < 1e-6) break;
        lambda = new_lambda;
    }
    return lambda;
}

double norm_2(const std::vector<double>& v) {
    double sum = 0.0;
    for (double val : v) sum += val * val;
    return sqrt(sum);
}

void test_correctness() {
    std::cout << "--- Test begin ---" << '\n';
    std::vector<int> test_sizes = {10, 50, 100, 500, 1000};
    
    for (int n : test_sizes) {
        CRSMatrix A = generate_spd_crs(n, 10);
        
        std::vector<double> x_true(n);
        std::mt19937 gen(123);
        std::uniform_real_distribution<double> dist_x(-1.0, 1.0);
        for (int i = 0; i < n; i++) x_true[i] = dist_x(gen);
        
        std::vector<double> b(n);
        multiply_CRS_vector(A, x_true, b);
        
        std::vector<double> x(n, 0.0);
        int iters = 0;
        
        SLE_Solver_CRS(A, b.data(), 1e-10, 10000, x.data(), iters);
        
        std::vector<double> Ax(n);
        multiply_CRS_vector(A, x, Ax);
        
        std::vector<double> residual(n);
        for (int i = 0; i < n; i++) residual[i] = Ax[i] - b[i];
        
        double norm_res = norm_2(residual);
        double norm_A = calc_matrix_2norm(A);
        double rel_error = norm_res / norm_A;
        
        std::cout << "n = " <<std::setw(5) << n 
             << " | Num inters: " << std::setw(4) << iters 
             << " | Rel. error: " << std::scientific << rel_error;
             
        if (rel_error < 1e-8) std::cout << " [PASSED]" << '\n';
        else std::cout << " [FAILED]" << '\n';
    }
    std::cout << '\n';
}

void test_performance() {
    std::cout << "--- Perf test begin ---" << '\n';
    int n = 100000; 
    int max_nz_per_row = 100;
    std::cout << " mat size " << n << "x" << n << "..." << std::flush;
    
    CRSMatrix A = generate_spd_crs(n, max_nz_per_row);
    std::cout << " A.nz: " << A.nz << '\n';
    
    std::vector<double> x_true(n, 1.0);
    std::vector<double> b(n);
    multiply_CRS_vector(A, x_true, b);
    std::vector<double> x(n, 0.0);
    
    std::vector<int> threads = {1, 2, 3, 4, 5, 6};
    double time_seq = 0.0;
    
    std::cout << std::fixed << std::setprecision(4);
    
    for (int p : threads) {
        omp_set_num_threads(p);
        fill(x.begin(), x.end(), 0.0);
        int iters = 0;
        auto start = std::chrono::steady_clock::now();
        SLE_Solver_CRS(A, b.data(), 1e-8, 1000, x.data(), iters);
        double time_spent = std::chrono::duration_cast<std::chrono::milliseconds>(std::chrono::steady_clock::now() - start).count() / 1000.l;

        if (p == 1) {
            time_seq = time_spent;
        }      
        double speedup = time_seq / time_spent;
        double efficiency = (speedup / p) * 100.0;

        std::cout << std::setw(10) << p << std::setw(15) << time_spent 
            << std::setw(15) << speedup << std::setw(19) << efficiency << "%" << '\n';
        
    }
}

int main() {
    test_correctness();
    test_performance();
    return 0;
}
\end{lstlisting}
\end{document}